%% toc: P2 23 23.1,23.2,23.3,23.4
\section{Proofs for the \(\Sigma^0_1\) definition of computably enumerable sets}
\label{sec:annex-p02-w02b-p02-034}

Notation is that of
\cref{def:p02-w02b-p02-034-machine,def:p02-w02b-p02-034-code,def:p02-w02b-p02-034-formulas}.
Four facts are used repeatedly.

\emph{(P1) Numerals and assignments.}  Induction on \(n\) shows that the
numeral \(\bar n=S^n0\) denotes \(n\) in \(\mathcal N\) under every
assignment.  Numerals are closed terms, so substituting several of them one
after another is the same as substituting them simultaneously, and
\cref{thm:fol-substitution} gives: \(\mathcal N\entails\varphi(\bar a_1,\ldots,\bar a_k)\)
iff \(\varphi\) holds in \(\mathcal N\) under an assignment sending
\(x_i\) to \(a_i\) for each \(i\).  The proofs pass between the two forms
without further comment.

\emph{(P2) Uniqueness of division.}  Let \(c\geq1\).  If \(cu+e=cu'+e'\) with
\(e,e'<c\), then \(u=u'\) and \(e=e'\).  Indeed, if \(u<u'\), then
\(cu+e<cu+c=c(u+1)\leq cu'\leq cu'+e'\), a contradiction; symmetrically
\(u'<u\) is impossible, so \(u=u'\), and cancelling \(cu\) gives \(e=e'\).

\emph{(P3) The beta formula.}  By \cref{lem:goedel-beta}(a), for all
\(a,i,c\in\N\), \(\mathcal N\entails\mathsf B(\bar a,\bar i,\bar c)\) iff
\(c=\beta(a,i)\).  By \cref{lem:goedel-beta}(b), for every finite sequence
\(c_0,\ldots,c_N\) there is \(a\) with \(\beta(a,i)=c_i\) for all
\(i\leq N\).  Since \(\bar i\) and \(S\bar i\) denote \(i\) and \(i+1\), the
instances \(\mathsf B(w,j,v)\) and \(\mathsf B(w,Sj,v)\) under an assignment
with \(j\mapsto i\) express \(v=\beta(w,i)\) and \(v=\beta(w,i+1)\).

\emph{(P4) Splitting off a digit.}  If \(z:\Z\to\{0,\ldots,b-1\}\) has
finitely many nonzero values and \(h\in\Z\), then reindexing the finite sum
gives
\(\sum_{j\in\N}z(h+j)b^j=z(h)+b\sum_{j\in\N}z(h+1+j)b^j\) and
\(\sum_{j\in\N}z(h-1-j)b^j=z(h-1)+b\sum_{j\in\N}z(h-2-j)b^j\).
The digit function of a configuration, \(i\mapsto\operatorname{dg}(T(i))\),
is such a \(z\), since \(\operatorname{dg}\) takes values below \(m=b\) and
\(\operatorname{dg}(\sqcup)=0\).

For a configuration \(C=(q,T,h)\) write
\(u_C\defeq\sum_{j\in\N}\operatorname{dg}(T(h+1+j))b^j\) and
\(\lambda_C\defeq\sum_{j\in\N}\operatorname{dg}(T(h-2-j))b^j\).  By (P4),
\[
 \operatorname{right}(C)=\operatorname{dg}(T(h))+b\,u_C,
 \qquad
 \operatorname{left}(C)=\operatorname{dg}(T(h-1))+b\,\lambda_C.
 \tag{E}
\]

\begin{annexitem}{annex:p02-w02b-p02-034-step}{thm:p02-w02b-p02-034-step}{One step of the machine as arithmetic}
\begin{proof}
Let \(C=(q,T,h)\).

\emph{(1)}  Suppose \(C\longrightarrow_{\mathcal T}C'\).  By the definition
of the one-step relation, \(q\) is not a halting state; let \(a=T(h)\) and
\(\delta(q,a)=(q',c,d)\), so \((q,a)\in D\), and \(C'=(q',T',h')\), where
\(T'\) agrees with \(T\) except that \(T'(h)=c\), and \(h'=h+1\) if \(d=R\),
\(h'=h-1\) if \(d=L\).  Interpret the variables of \(\Delta_{q,a}\) by
\(p=\operatorname{st}(q)\), \(l=\operatorname{left}(C)\),
\(r=\operatorname{right}(C)\), \((p',l',r')=\operatorname{enc}(C')\), and
the witness \(u=u_C\).  The conjunct \(p=\overline{\operatorname{st}(q)}\)
holds, \(r=(\bar b\times u)+\overline{\operatorname{dg}(a)}\) holds by (E),
and \(p'=\overline{\operatorname{st}(q')}\) holds because \(q'\) is the state
of \(C'\).

\emph{Case \(d=R\).}  Every cell \(h+1+j\) differs from \(h\), so
\(\operatorname{right}(C')=\sum_j\operatorname{dg}(T'(h+1+j))b^j=u_C\), which
is \(r'=u\).  By (P4) applied to \(T'\) at \(h'-1=h\),
\[
 \operatorname{left}(C')=\sum_j\operatorname{dg}(T'(h-j))b^j
 =\operatorname{dg}(T'(h))+b\sum_j\operatorname{dg}(T'(h-1-j))b^j
 =\operatorname{dg}(c)+b\operatorname{left}(C),
\]
because the cells \(h-1-j\) differ from \(h\).  This is
\(l'=(\bar b\times l)+\overline{\operatorname{dg}(c)}\).

\emph{Case \(d=L\).}  Put \(g=\operatorname{dg}(T(h-1))\), so \(g<b\).
The cells \(h-2-j\) differ from \(h\), hence
\(\operatorname{left}(C')=\sum_j\operatorname{dg}(T'(h'-1-j))b^j=\lambda_C\),
and (E) gives \(\operatorname{left}(C)=b\operatorname{left}(C')+g\), which is
\(l=(\bar b\times l')+\bar g\).  Applying (P4) twice to \(T'\) at \(h'=h-1\),
\[
 \operatorname{right}(C')
 =\operatorname{dg}(T'(h-1))+b\bigl(\operatorname{dg}(T'(h))+b\sum_j\operatorname{dg}(T'(h+1+j))b^j\bigr)
 =g+b\bigl(\operatorname{dg}(c)+b\,u_C\bigr),
\]
which is the second conjunct of the disjunct for this \(g\).  So
\(\Delta_{q,a}\) holds in both cases, hence \(\exists u\bigvee_D\Delta\) holds,
and by (P1)
\(\mathcal N\entails\mathsf{Step}_{\mathcal T}(\overline{\operatorname{enc}(C)},\overline{\operatorname{enc}(C')})\).

\emph{(2)}  Suppose
\(\mathcal N\entails\mathsf{Step}_{\mathcal T}(\overline{\operatorname{enc}(C)},\bar p',\bar l',\bar r')\).
Then there are \(u\in\N\) and \((\hat q,\hat a)\in D\) such that
\(\Delta_{\hat q,\hat a}\) holds with \(p=\operatorname{st}(q)\),
\(l=\operatorname{left}(C)\), \(r=\operatorname{right}(C)\).  Its first
conjunct gives \(\operatorname{st}(q)=\operatorname{st}(\hat q)\), so
\(q=\hat q\) by injectivity of \(\operatorname{st}\); thus \(q\) is not
halting.  Its second conjunct gives
\(\operatorname{right}(C)=bu+\operatorname{dg}(\hat a)\) with
\(\operatorname{dg}(\hat a)<b\), while (E) gives
\(\operatorname{right}(C)=b\,u_C+\operatorname{dg}(T(h))\) with
\(\operatorname{dg}(T(h))<b\).  By (P2), \(u=u_C\) and
\(\operatorname{dg}(\hat a)=\operatorname{dg}(T(h))\), so \(\hat a=T(h)\) by
injectivity of \(\operatorname{dg}\).  Hence \(\delta(q,T(h))=\delta(\hat q,\hat a)=(q',c,d)\)
is defined, \(C\) has the successor \(C'=(q',T',h')\) described in (1), and
\(p'=\operatorname{st}(q')\).

If \(d=R\), the remaining conjuncts give
\(l'=b\operatorname{left}(C)+\operatorname{dg}(c)\) and \(r'=u=u_C\); by the
computation in (1) these are \(\operatorname{left}(C')\) and
\(\operatorname{right}(C')\).  If \(d=L\), some \(g<b\) satisfies
\(\operatorname{left}(C)=bl'+g\) and
\(r'=g+b(\operatorname{dg}(c)+bu)\).  By (1),
\(\operatorname{left}(C)=b\operatorname{left}(C')+\operatorname{dg}(T(h-1))\)
with \(\operatorname{dg}(T(h-1))<b\), so (P2) gives
\(l'=\operatorname{left}(C')\) and \(g=\operatorname{dg}(T(h-1))\); then, with
\(u=u_C\), the computation in (1) gives \(r'=\operatorname{right}(C')\).  In
both cases \((p',l',r')=\operatorname{enc}(C')\).
\end{proof}
\end{annexitem}

\begin{annexitem}{annex:p02-w02b-p02-034-init}{thm:p02-w02b-p02-034-init}{The initial configuration as arithmetic}
\begin{proof}
Write \(w=\operatorname{bin}(x)\), a nonempty word over
\(\{\mathtt 0,\mathtt 1\}\) (\(\operatorname{bin}(0)\) has one symbol, and
the base-two expansion of a positive number has at least one digit).  The
tape of \(C^x_0\) is \(T_w\), with \(T_w(i)=w(i)\) for \(0\le i<|w|\) and
\(T_w(i)=\sqcup\) otherwise, and the head is at \(0\).  Every cell \(-1-j\)
is blank, so \(\operatorname{left}(C^x_0)=0\), and
\[
 \operatorname{right}(C^x_0)=\sum_{j\in\N}\operatorname{dg}(T_w(j))b^j
 =\sum_{j<|w|}\operatorname{dg}(w(j))b^j,
 \tag{R}
\]
because \(\operatorname{dg}(\sqcup)=0\).  For a word
\(a_0\cdots a_k\) over \(\{\mathtt 0,\mathtt 1\}\) let \(e_i\in\{0,1\}\) be the
number with \(a_i=\mathtt 0\) iff \(e_i=0\), and call
\(\sum_{i\le k}e_i2^{k-i}\) its \emph{value}.

\emph{Uniqueness of binary words (U).}  If two words \(a_0\cdots a_k\) and
\(a'_0\cdots a'_{k'}\) with \(a_0=a'_0=\mathtt 1\) have the same value, they
are equal.  The value \(V\) of the first satisfies
\(2^k\le V\le\sum_{i\le k}2^{k-i}=2^{k+1}-1\), and likewise for the second
with \(k'\); these ranges are disjoint for different lengths, so \(k=k'\).
Induct on \(k\).  For \(k=0\) both words are \(\mathtt 1\).  For \(k\geq1\),
\(V=2V^-+e_k\), where \(V^-\) is the value of \(a_0\cdots a_{k-1}\), and
likewise \(V=2V'^-+e'_k\); by (P2) with \(c=2\), \(e_k=e'_k\) and
\(V^-=V'^-\), and the two prefixes, both beginning with \(\mathtt 1\), are
equal by the induction hypothesis.

\emph{The recurrences.}  Let \(a_0\cdots a_k\) be a word over
\(\{\mathtt 0,\mathtt 1\}\) with digits \(e_i\) as above, and suppose numbers
\(v_j,t_j,\pi_j\) (\(j\le k+1\)) satisfy \(v_0=t_0=0\), \(\pi_0=1\), and for
\(j\le k\): \(\pi_{j+1}=b\pi_j\), \(v_{j+1}=2v_j+e_j\) and
\(t_{j+1}=t_j+\operatorname{dg}(a_j)\pi_j\).  Induction on \(j\le k+1\) gives
\[
 \pi_j=b^j,\qquad v_j=\sum_{i<j}e_i2^{j-1-i},\qquad
 t_j=\sum_{i<j}\operatorname{dg}(a_i)b^i:
 \tag{C}
\]
at \(j=0\) the sums are empty; and if (C) holds at \(j\le k\), then
\(\pi_{j+1}=b\cdot b^j\),
\(v_{j+1}=2\sum_{i<j}e_i2^{j-1-i}+e_j=\sum_{i<j+1}e_i2^{j-i}\), and
\(t_{j+1}=\sum_{i<j}\operatorname{dg}(a_i)b^i+\operatorname{dg}(a_j)b^j\).
Conversely, the numbers defined by (C) satisfy the recurrences, by the same
computation.  In particular \(v_{k+1}\) is the value of the word and
\(t_{k+1}=\sum_{i\le k}\operatorname{dg}(a_i)b^i\).

\emph{Left to right.}  Suppose
\(\mathcal N\entails\mathsf{Init}_{\mathcal T}(\bar x,\bar r)\), with witnesses
\(k,\sigma_1,\sigma_2,\sigma_3,\sigma_4\) for \(k,s_1,s_2,s_3,s_4\).  Put
\(e_j=\beta(\sigma_1,j)\), \(v_j=\beta(\sigma_2,j)\), \(t_j=\beta(\sigma_3,j)\)
and \(\pi_j=\beta(\sigma_4,j)\).  By (P3) the first line gives
\(v_0=t_0=0\) and \(\pi_0=1\).  The marked bounded quantifier
\(\forall j<Sk\) ranges over \(j\le k\); for each such \(j\) its witnesses
\(e,v,v',t,t',\pi,\pi'\) are, by (P3), \(e_j,v_j,v_{j+1},t_j,t_{j+1},\pi_j,\pi_{j+1}\),
and the matrix gives \(\pi_{j+1}=b\pi_j\) together with either
\(e_j=0\), \(v_{j+1}=2v_j\), \(t_{j+1}=t_j+\operatorname{dg}(\mathtt 0)\pi_j\),
or \(e_j=1\), \(v_{j+1}=2v_j+1\), \(t_{j+1}=t_j+\operatorname{dg}(\mathtt 1)\pi_j\).
Let \(a_j\) be \(\mathtt 0\) if \(e_j=0\) and \(\mathtt 1\) if \(e_j=1\).  The
recurrences hold for the word \(a_0\cdots a_k\), so by (C) and the last line,
\(x=v_{k+1}\) is its value and \(r=t_{k+1}=\sum_{i\le k}\operatorname{dg}(a_i)b^i\);
moreover \(k=0\) or \(e_0=1\).

The word \(a_0\cdots a_k\) is \(\operatorname{bin}(x)\).  If \(k=0\), it is
\(a_0\) and \(x=e_0\); \(\operatorname{bin}(0)=\mathtt 0\), and
\(\operatorname{bin}(1)=\mathtt 1\) because the base-two expansion of \(1\)
without leading zeros is the single digit \(1\).  If \(k\ge1\), then
\(e_0=1\), so \(x\ge2^k>0\), and \(\operatorname{bin}(x)\) is the base-two
expansion of \(x\) without leading zeros, a word beginning with
\(\mathtt 1\) whose value is \(x\); by (U) it equals \(a_0\cdots a_k\).  By (R),
\(r=\operatorname{right}(C^x_0)\).

\emph{Right to left.}  Let \(r=\operatorname{right}(C^x_0)\) and write
\(\operatorname{bin}(x)=a_0\cdots a_k\) with digits \(e_i\).  Define
\(v_j,t_j,\pi_j\) for \(j\le k+1\) by (C).  By (P3) there are
\(\sigma_1,\ldots,\sigma_4\) with \(\beta(\sigma_1,j)=e_j\) for \(j\le k\) and
\(\beta(\sigma_2,j)=v_j\), \(\beta(\sigma_3,j)=t_j\), \(\beta(\sigma_4,j)=\pi_j\)
for \(j\le k+1\).  Interpret \(k,s_1,\ldots,s_4\) by
\(k,\sigma_1,\ldots,\sigma_4\).  The first line holds by (P3).  For each
\(j\le k\), the witnesses \(e_j,v_j,v_{j+1},t_j,t_{j+1},\pi_j,\pi_{j+1}\)
make the matrix true: the \(\mathsf B\)-conjuncts by (P3), and the arithmetic
conjuncts because the numbers (C) satisfy the recurrences, with the disjunct
chosen by \(e_j\).  The value \(v_{k+1}\) of \(\operatorname{bin}(x)\) is
\(x\), by the definition of the base-two expansion (and
\(\operatorname{bin}(0)=\mathtt 0\) has value \(0\)); and
\(t_{k+1}=r\) by (R).  Finally, if \(k\ge1\) then \(x\ne0\), because
\(\operatorname{bin}(0)\) has one symbol, so \(\operatorname{bin}(x)\) has no
leading zero and \(e_0=1\).  Hence
\(\mathcal N\entails\mathsf{Init}_{\mathcal T}(\bar x,\bar r)\).
\end{proof}
\end{annexitem}

\begin{annexitem}{annex:p02-w02b-p02-034-halt}{thm:p02-w02b-p02-034-halt}{Halting is \(\Sigma^0_1\)-definable}
\begin{proof}
\emph{Classification.}  The formula \(\mathsf B\) of \cref{lem:goedel-beta}
is built from atomic formulas by \(\land\) and marked bounded existential
quantifiers, with no negation sign and no implication; so it, and every
substitution instance of it with renamed bound variables, is a bounded
formula in negation normal form, that is, a \(\Sigma^0_0\) formula
(\cref{sec:decision-hierarchy-02}).  Atomic formulas and \(\bottom\) are
\(\Sigma^0_0\) as well.  By the definition of the syntactic hierarchy, the
\(\Sigma^0_1\) formulas in negation normal form contain every
\(\Sigma^0_0\) formula and are closed under \(\land\), \(\lor\), both marked
bounded quantifiers and unbounded \(\exists\).  Hence:
each \(\Delta_{q,a}\), and the disjunction over \(D\), is built from atoms by
\(\land,\lor\), so it is \(\Sigma^0_0\) and \(\mathsf{Step}_{\mathcal T}\),
an unbounded \(\exists\) of it, is \(\Sigma^0_1\).  In
\(\mathsf{Init}_{\mathcal T}\) the matrix under \(\forall j<Sk\) is an
unbounded existential block over a conjunction of \(\mathsf B\)-instances,
atoms and a disjunction of conjunctions of atoms, so it is \(\Sigma^0_1\);
since \(j\) does not occur in the bound \(Sk\), \(\forall j<Sk\) is a marked
bounded quantifier, and the result is \(\Sigma^0_1\); the remaining
conjunctions, the disjunction \(k=\bar0\lor\mathsf B(s_1,\bar0,\bar1)\) and the
outer existential block keep it \(\Sigma^0_1\).  In
\(\mathsf{Halt}_{\mathcal T}\) the same reasoning applies: the matrix under
\(\forall k<n\) (with \(k\) not occurring in \(n\)) is an existential block
over \(\mathsf B\)-instances and the \(\Sigma^0_1\) formula
\(\mathsf{Step}_{\mathcal T}\); the last conjunct is an existential
quantifier over \(\Sigma^0_0\) material; \(\mathsf{Init}_{\mathcal T}\) is
\(\Sigma^0_1\); and conjunction and the outer existential block preserve the
class.  The formula contains no negation sign and no implication, so its
fixed negation-normal-form conversion makes no rewrite and at most renames
bound variables, which changes no generating step.  So
\(\mathsf{Halt}_{\mathcal T}\) is a \(\Sigma^0_1\) formula.

\emph{From truth to halting.}  Suppose
\(\mathcal N\entails\mathsf{Halt}_{\mathcal T}(\bar x)\), with witnesses
\(n,\alpha_1,\alpha_2,\alpha_3,\rho_0\) for \(n,w_1,w_2,w_3,r_0\).  Put
\(P_k=\beta(\alpha_1,k)\), \(L_k=\beta(\alpha_2,k)\), \(R_k=\beta(\alpha_3,k)\).
By \cref{thm:p02-w02b-p02-034-init}, \(\rho_0=\operatorname{right}(C^x_0)\)
and \(\operatorname{left}(C^x_0)=0\); by (P3) the next three conjuncts give
\(P_0=\operatorname{st}(q_0)\), \(L_0=0\), \(R_0=\rho_0\).  So
\((P_0,L_0,R_0)=\operatorname{enc}(C^x_0)\).  For each \(k<n\) the instance
\(k\) of \(\forall k<n\) supplies witnesses which by (P3) are
\(P_k,L_k,R_k,P_{k+1},L_{k+1},R_{k+1}\), so
\[
 \mathcal N\entails\mathsf{Step}_{\mathcal T}(\bar P_k,\bar L_k,\bar R_k,\bar P_{k+1},\bar L_{k+1},\bar R_{k+1})
 \qquad(k<n).
 \tag{S}
\]
Claim: for each \(k\le n\), the computation of \(\mathcal T\) on
\(\operatorname{bin}(x)\) has an entry \(C_k\) at position \(k\), and
\(\operatorname{enc}(C_k)=(P_k,L_k,R_k)\).  For \(k=0\), \(C_0=C^x_0\).  If
the claim holds for \(k<n\), then (S) and
\cref{thm:p02-w02b-p02-034-step}(2) show that \(C_k\) is not halting and has
a successor \(C'\) with \(\operatorname{enc}(C')=(P_{k+1},L_{k+1},R_{k+1})\).
The computation is the maximal sequence that starts at \(C^x_0\) and follows
\(\longrightarrow_{\mathcal T}\); a finite sequence ending at \(C_k\) could
be extended by \(C'\), so the computation has an entry \(C_{k+1}\) at
position \(k+1\) with \(C_k\longrightarrow_{\mathcal T}C_{k+1}\).  The
one-step relation is functional, since the successor is determined by
\(\delta(q,T(h))\) through the displayed \(T'\) and \(h'\), so \(C_{k+1}=C'\).
This proves the claim.  The last conjunct and (P3) give
\(P_n\in\{\operatorname{st}(q_{\mathrm{acc}}),\operatorname{st}(q_{\mathrm{rej}})\}\),
and \(P_n\) is the state number of \(C_n\); by injectivity of
\(\operatorname{st}\), \(C_n\) is in a halting state.  So the computation
halts.

\emph{From halting to truth.}  Suppose the computation halts: it is a finite
sequence \(C_0,\ldots,C_n\) with \(C_0=C^x_0\),
\(C_k\longrightarrow_{\mathcal T}C_{k+1}\) for \(k<n\), and \(C_n\) in a
halting state.  By (P3) choose \(\alpha_1,\alpha_2,\alpha_3\) with
\((\beta(\alpha_1,k),\beta(\alpha_2,k),\beta(\alpha_3,k))=\operatorname{enc}(C_k)\)
for every \(k\le n\), and put \(\rho_0=\operatorname{right}(C^x_0)\).
Interpret \(n,w_1,w_2,w_3,r_0\) by \(n,\alpha_1,\alpha_2,\alpha_3,\rho_0\).
Then \(\mathsf{Init}_{\mathcal T}(x,r_0)\) holds by
\cref{thm:p02-w02b-p02-034-init}; the three \(\mathsf B\)-conjuncts at
\(\bar0\) hold by (P3), since \(C_0\) has state \(q_0\), left code \(0\) and
right code \(\rho_0\).  For \(k<n\), the witnesses
\(\operatorname{enc}(C_k)\) and \(\operatorname{enc}(C_{k+1})\) make the
\(\mathsf B\)-conjuncts true by (P3) and \(\mathsf{Step}_{\mathcal T}\) true by
\cref{thm:p02-w02b-p02-034-step}(1).  The last conjunct holds with
\(p_\ast=\beta(\alpha_1,n)\), the state number of \(C_n\), which is
\(\operatorname{st}(q_{\mathrm{acc}})\) or \(\operatorname{st}(q_{\mathrm{rej}})\).
Hence \(\mathcal N\entails\mathsf{Halt}_{\mathcal T}(\bar x)\).
\end{proof}
\end{annexitem}

\begin{annexitem}{annex:p02-w02b-p02-034-ce-sigma-one}{thm:p02-w02b-p02-034-ce-sigma-one}{Every computably enumerable set is \(\Sigma^0_1\)-definable}
\begin{proof}
Let \(A\subseteq\N\) be computably enumerable.  By the definition in
\cref{sec:turing-machines-03}, \(A=\operatorname{dom}(f)\) for a partial
computable \(f:\N\rightharpoonup\N\), and by \cref{def:partial-computable}
some Turing machine \(\mathcal T\) satisfies, for every \(x\in\N\): if
\(f(x)\) is defined, the computation of \(\mathcal T\) on
\(\operatorname{bin}(\operatorname{code}_1(x))=\operatorname{bin}(x)\) reaches
\(q_{\mathrm{acc}}\); if \(f(x)\) is undefined, that computation is infinite.
These inputs are words over \(\{\mathtt 0,\mathtt 1\}\), and a machine used to
compute a numerical function has \(\mathtt 0,\mathtt 1\) in its input alphabet
(\cref{sec:turing-machines-01}), so \cref{def:p02-w02b-p02-034-machine}
applies to \(\mathcal T\).  A computation that reaches
\(q_{\mathrm{acc}}\) halts, and an infinite one does not.  Hence \(x\in A\)
iff the computation on \(\operatorname{bin}(x)\) halts, iff
\(\mathcal N\entails\mathsf{Halt}_{\mathcal T}(\bar x)\) by
\cref{thm:p02-w02b-p02-034-halt}.  Since \(\mathsf{Halt}_{\mathcal T}(x)\) is a
\(\Sigma^0_1\) formula with free variable \(x\) only, \(A\) is
\(\Sigma^0_1\)-definable, that is, \(A\in\boldsymbol\Sigma^0_1\).
\end{proof}
\end{annexitem}

\begin{annexitem}{annex:p02-w02b-p02-034-ta-not-ce}{thm:p02-w02b-p02-034-ta-not-ce}{True arithmetic is not computably enumerable}
\begin{proof}
Suppose \(\mathrm{TA}\) were computably enumerable.  By
\cref{thm:p02-w02b-p02-034-ce-sigma-one} there is a \(\Sigma^0_1\) formula
\(\mathsf{Tr}(x)\) such that, for every \(c\in\N\), \(c\in\mathrm{TA}\) iff
\(\mathcal N\entails\mathsf{Tr}(\bar c)\).  Let \(\sigma\) be any
\(\Lang_A\)-sentence.  The code \(\ulcorner\sigma\urcorner\) in
\cref{def:p02-w02-p02-040-so-coding} is the code of
\cref{sec:arithmetic-coding-04}, the one used in the definition of
arithmetical truth definitions (\cref{sec:second-incompleteness-04}).  By the
definition of \(\mathrm{TA}\),
\(\mathcal N\entails\mathsf{Tr}(\overline{\ulcorner\sigma\urcorner})\) iff
\(\ulcorner\sigma\urcorner\in\mathrm{TA}\) iff \(\mathcal N\entails\sigma\).
Since \(\mathcal N\entails\psi\leftrightarrow\chi\) holds for sentences
\(\psi,\chi\) exactly when both are true or both are false,
\(\mathcal N\entails\mathsf{Tr}(\overline{\ulcorner\sigma\urcorner})\leftrightarrow\sigma\)
for every sentence \(\sigma\).  So \(\mathsf{Tr}\) is an arithmetical truth
definition, which \cref{thm:tarski-undefinability} excludes.  Hence
\(\mathrm{TA}\) is not computably enumerable.
\end{proof}
\end{annexitem}
